\scanpage{34}
mais on fera attention que l'inclusion % L1
\[
\begin{aligned}
\deleted{\mathbb Z}\mathbb Q(-1)\otimes\mathcal M^{+i}(X)&\hookrightarrow\mathcal M^{+(i+2)}(X)\qquad\text{plus gén.}\\ % L2
\deleted{\mathbb Z}\mathbb Q(-j)\otimes\mathcal M^{+i}(X)&\unsure{\lhook\joinrel-|\hookrightarrow}M^{+(i+2j)}(X) % L3
\end{aligned}
\]
est \uline{stricte} en général, i.e. n'est pas % L4
une équivalence de catégories. % L5

Noter que $\mathcal M^{+i}(X)$ est \uline{épaisse} dans $\mathcal M^{+j}(X)$ % L6
% slanted
\margin{indices en bas} % L6a
\deleted{[elle est définie, dans la catégorie}, et de % L7
$=$ $\mathcal M^i(X)$ épaisse dans $\mathcal M^j(X)$ $(i\leq j)$. % L8

Les quotients $G_i^+(X)=\mathrm{Gr}^{\unsure{F|\mathscr F}}_i(\mathcal M^+(X)\deleted{\unsure{)_{i-1}}})\inserted{\simeq\mathcal M_i^+(X)/\mathcal M_{i-1}^+(X)}$ et $G_i(X)=\mathrm{Gr}^{F}_i(\mathcal M(X))$ % L9
$\simeq\mathcal M_i(X)/\mathcal M_{i-1}(X)$ sont fort intéressants. % L10

Noter que les \deleted{\unsure{M}}$G_i(X)$ sont donc \deleted{\unsure{isomorphes}} \inserted{équivalentes} % L11
par twisting $G_i(X)\overset{\mathbb Z(-j)}{\simeq}G_{i+2j}(X)$, en % L12
particulier tous \deleted{\unsure{isomorphes}} \inserted{équiv. can.} à $G_0(X)$. % L13
D'ailleurs $\unsure{G_i^+|G\mathrm r_i^+|\mathcal G_i^+}(X)\underset{\sim}{\hookrightarrow}G_i(X)\overset{\sim}{\simeq}G_0(X)$ % L14
\inserted{équiv. à une} sous catégorie pleine et \uline{épaisse} de $G_i(X)\overset{\sim}{\simeq} % L15
G_0(X)$. % L16

% archive sheet number: 33
